\documentclass[11pt,a4paper]{article}
\usepackage[utf8]{inputenc}
\usepackage[T1]{fontenc}
\usepackage[margin=1in]{geometry}
\usepackage{hyperref}
\usepackage{xcolor}
\usepackage{listings}
\usepackage{upquote}
\usepackage{enumitem}
\usepackage{tcolorbox}
\usepackage{booktabs}
\usepackage{amssymb}
\tcbuselibrary{skins}

\lstset{
    basicstyle=\ttfamily\small,
    breaklines=true,
    frame=single,
    numbers=left,
    numberstyle=\tiny,
    backgroundcolor=\color{gray!10},
    keywordstyle=\color{blue},
    commentstyle=\color{green!50!black},
    stringstyle=\color{red},
    showstringspaces=false
}

\newtcolorbox{notebox}{colback=blue!5,colframe=blue!40!black,title=Note,fonttitle=\bfseries}
\newtcolorbox{bonusbox}{colback=green!5,colframe=green!40!black,title=Bonus Opportunity,fonttitle=\bfseries}

\title{\textbf{Final Project --- Stretch Robot Proof of Concept}}
\author{
    Robotics and Cyber-Physical Systems \\
    School of Sciences and Engineering \\
    Tecnol\'ogico de Monterrey
}
\date{}

\begin{document}
\maketitle

\section*{Overview}

For the final project, your team will identify a practical use case for the
Hello Robot Stretch and implement a proof-of-concept demonstration. Your demo
should provide convincing evidence that the idea is feasible with the robot's
hardware, sensors, and workspace.

The goal is not a finished production system. A successful project has a
focused objective, demonstrates meaningful robot behavior, and communicates
its limitations honestly.

Use the \texttt{stretch\_mujoco\_digital\_twin} simulator as the primary
development platform. The current API is organized as follows:

\begin{itemize}
    \item The simulated \texttt{Robot} mirrors the important parts of the
          physical \texttt{stretch\_body.robot} interface.
    \item \texttt{stretch\_tools} provides environment-agnostic tools such as
          normalized velocity control, teleoperation, cameras, state control,
          transforms, and LiDAR helpers.
    \item \texttt{sim\_config.py} configures simulator-only details such as
          objects, RoboCasa, and initial poses.
\end{itemize}

Keep environment-specific setup separate from the main behavior whenever
possible. This makes the core script easier to transfer to the physical robot.

\section*{Choosing a Use Case}

Any application that makes genuine use of Stretch's capabilities is
acceptable. Possible directions include:

\begin{itemize}
    \item Assistive robotics
    \item Object detection, manipulation, or sorting
    \item Inspection and environmental monitoring
    \item Human--robot interaction
    \item Autonomous navigation or fetch tasks
\end{itemize}

Choose a specific, testable objective. ``Retrieve a red object from a table''
is easier to evaluate than ``build a general-purpose home robot.'' Your scope
should be ambitious enough to be interesting, but small enough to demonstrate
reliably within the project timeline.

\section*{Technical Expectations}

The project should combine multiple ideas from the course and include some
autonomous behavior. Pure teleoperation is not sufficient by itself.

Depending on the use case, relevant components may include:

\begin{itemize}
    \item Position and velocity commands through the \texttt{Robot} API
    \item \texttt{NormVelController}, \texttt{StateController}, or
          \texttt{TeleopProvider}
    \item RGB or depth-camera perception and visual servoing
    \item Object localization with \texttt{RobotTransforms}
    \item LiDAR-based navigation or obstacle avoidance
    \item Simulator objects or RoboCasa configured through
          \texttt{sim\_config.py}
\end{itemize}

You are not required to use every feature. Select the tools that support your
application and explain your design choices.

\begin{notebox}
Test incrementally. Keep speeds conservative, enable collision management when
running on the physical robot, and provide a reliable way to stop autonomous
motion. A complex demo is not successful if it cannot be operated safely and
repeatably.
\end{notebox}

\section*{Recommended Project Structure}

Start from \texttt{stretch\_boilerplate.py} or
\texttt{teleop\_boilerplate.py}. Use conditional imports so the main script
can run in either environment when practical:

\begin{lstlisting}[language=Python]
try:
    import stretch_body.robot as robot
except ImportError:
    import stretch_mujoco_api.robot as robot

from stretch_tools import NormVelController
\end{lstlisting}

Simulator-only scene choices belong in \texttt{sim\_config.py}, not in the
portable behavior script. If your project needs custom meshes or textures,
include those assets and the corresponding configuration with your submission.

\section*{Deliverables}

Every team must submit:

\begin{enumerate}
    \item \textbf{Source code} --- A working script or small set of scripts
          that runs in the simulator. Include any required
          \texttt{sim\_config.py}, custom assets, and brief comments for
          non-obvious logic. Remove temporary diagnostics and unused code.

    \item \textbf{Written report} (2--3 pages) --- A concise document covering:
    \begin{itemize}
        \item The use case and why Stretch is suitable
        \item The system architecture and control/perception approach
        \item How the system was tested and what the results show
        \item Known limitations, safety considerations, and next steps
        \item Any differences between simulated and physical operation
    \end{itemize}

    \item \textbf{Video} --- A recording that clearly shows the complete demo
          running. It may be a screen capture or an edited/narrated video, but
          it should be understandable without the class presentation.
\end{enumerate}

\textbf{File naming:}
\texttt{Team\_[TeamName]\_FinalProject.zip}. Include the source code, report,
video or video link, configuration, and required assets.

\section*{Presentation}

Each team will give a 10--15 minute presentation, including questions. It
should contain:

\begin{itemize}
    \item \textbf{Motivation} --- the problem and intended user
    \item \textbf{Approach} --- the main perception and control decisions
    \item \textbf{Demo} --- a live simulator demonstration, with the submitted
          video available as a fallback
    \item \textbf{Evaluation} --- what worked, what failed, and what you learned
    \item \textbf{Q\&A} --- answers that demonstrate understanding of your code
\end{itemize}

\begin{bonusbox}
Teams that successfully demonstrate their proof of concept on the
\textbf{physical Stretch robot} are exempt from the class presentation. All
three project deliverables are still required.

Arrange a supervised robot session using the instructor's scheduling tool.
One or two team representatives may attend, and multiple teams may coordinate
to share a session.
\end{bonusbox}

\section*{Grading}

The project is graded across four equally weighted areas:

\begin{center}
\begin{tabular}{llp{8.4cm}}
\toprule
\textbf{Area} & \textbf{Weight} & \textbf{What we are looking for} \\
\midrule
\textbf{Idea \& Feasibility} & 25\% &
    The use case fits Stretch, has a clear objective, and is scoped as a
    credible proof of concept. \\[6pt]
\textbf{Implementation} & 25\% &
    The demo works reliably, includes meaningful autonomous behavior, uses
    appropriate robot capabilities, and is implemented clearly and safely. \\[6pt]
\textbf{Presentation} & 25\% &
    The team explains the problem, approach, results, and limitations clearly;
    the demonstration supports the team's claims. \\[6pt]
\textbf{Video} & 25\% &
    The recording clearly shows the complete behavior and can be understood
    without additional explanation. \\
\bottomrule
\end{tabular}
\end{center}

\begin{notebox}
There is no single correct implementation. Grading will consider the ambition
and difficulty of the selected problem together with the quality and
reliability of the achieved result.
\end{notebox}

\section*{Submission Checklist}

Before submitting, verify that:

\begin{itemize}[label=$\square$]
    \item A fresh checkout can install the dependencies and run the demo.
    \item All required files and assets use repository-relative paths.
    \item The simulator configuration needed for the demo is included.
    \item The robot stops safely when the program exits or loses its target.
    \item Temporary debug output and unrelated files have been removed.
    \item The report and video explain the final behavior and its limitations.
\end{itemize}

\end{document}
